% =====================================================
% J. J. Thomson (1897): Nachweis des Elektrons
% =====================================================
\subsection{J. J. Thomson (1897): Nachweis des Elektrons}
\label{sec:thomson-nachweis-elektron}

Die Experimente mit Kathodenstrahlen\index{Kathodenstrahl} hatten gezeigt, dass in
Entladungsröhren\index{Entladungsröhre} etwas geschieht, das sich nicht einfach als
gewöhnliches Licht erklären ließ. Die Strahlen breiteten sich geradlinig aus, konnten
Schatten erzeugen und brachten Glaswände zum Leuchten. Zugleich ließen sie sich durch
elektrische und magnetische Felder\index{Magnetfeld} ablenken. Genau diese Ablenkbarkeit
war der entscheidende Hinweis: Kathodenstrahlen mussten mit elektrischer Ladung
verbunden sein.

J. J. Thomson\index{Thomson, J. J.} untersuchte diese Strahlen im Jahr 1897 besonders
systematisch. Seine zentrale Frage lautete: Handelt es sich bei Kathodenstrahlen um
eine Wellenerscheinung oder um einen Strom kleiner geladener Teilchen? Um diese Frage
zu klären, ließ er Kathodenstrahlen durch elektrische und magnetische Felder laufen
und beobachtete ihre Ablenkung.

Wenn ein Teilchen elektrisch geladen ist, erfährt es in einem elektrischen Feld
\index{elektrisches Feld} eine Kraft. In einem Magnetfeld erfährt es ebenfalls eine
Kraft, sofern es sich bewegt. Aus der Stärke der Ablenkung konnte Thomson schließen,
dass die Kathodenstrahlen aus negativ geladenen Teilchen bestehen mussten. Diese
Teilchen wurden später Elektronen\index{Elektron} genannt.

Besonders wichtig war, dass Thomson nicht nur eine qualitative Beobachtung machte.
Er bestimmte das Verhältnis von Ladung zu Masse\index{Ladung-Masse-Verhältnis} dieser
Teilchen. Dieses Verhältnis wird als
\[
\frac{e}{m}
\]
bezeichnet. Thomson fand, dass dieses Verhältnis sehr viel größer war als bei bekannten
Ionen\index{Ion}. Daraus folgte: Entweder war die Ladung dieser Teilchen ungewöhnlich
groß, oder ihre Masse war außerordentlich klein. Die spätere Deutung zeigte, dass vor
allem Letzteres zutraf: Das Elektron besitzt eine extrem kleine Masse.

Ein weiterer entscheidender Punkt war die Unabhängigkeit vom verwendeten Material.
Thomson stellte fest, dass das Ladung-Masse-Verhältnis der Kathodenstrahlteilchen nicht
wesentlich davon abhing, welches Gas sich in der Röhre befand oder aus welchem Material
die Kathode bestand. Das war ein starkes Argument dafür, dass es sich nicht um eine
besondere Eigenschaft eines einzelnen Stoffes handelte, sondern um einen allgemeinen
Bestandteil der Materie\index{Materie}.

Damit wurde eine alte Vorstellung erschüttert. Atome\index{Atom} galten lange Zeit als
unteilbare Grundbausteine der Natur. Wenn Kathodenstrahlen aber aus Teilchen bestanden,
die viel leichter waren als Atome und aus unterschiedlichen Stoffen hervortreten
konnten, dann mussten Atome selbst eine innere Struktur besitzen. Thomsons Experiment
war daher mehr als der Nachweis eines neuen Teilchens. Es war ein Wendepunkt im
Verständnis der Materie.

Thomson selbst sprach zunächst von \glqq Korpuskeln\grqq, also kleinen Teilchen.
Der Name Elektron setzte sich später durch. Rückblickend gilt sein Experiment von
1897 als entscheidender Nachweis des Elektrons als eigenständigem Bestandteil der
Materie. Aus den rätselhaften Kathodenstrahlen wurde damit der erste direkte Zugang
zu einem subatomaren Teilchen\index{subatomares Teilchen}.

\begin{PhysicsBox}[Thomsons entscheidende Idee]
	\label{box:thomson-entscheidende-idee}
	Thomson zeigte, dass Kathodenstrahlen aus negativ geladenen Teilchen bestehen.
	Dazu untersuchte er ihre Ablenkung in elektrischen und magnetischen Feldern.
	
	Entscheidend war das Verhältnis von Ladung zu Masse:
	\[
	\frac{e}{m}.
	\]
	
	Da dieses Verhältnis viel größer war als bei bekannten Ionen, musste das
	Teilchen eine sehr kleine Masse besitzen. Damit wurde klar: In der Materie
	gibt es Bestandteile, die kleiner sind als Atome.
\end{PhysicsBox}

\begin{HistoryBox}[Warum Thomson die Atomvorstellung veränderte]
	\label{box:thomson-atomvorstellung}
	Vor Thomson wurden Atome vielfach als unteilbare Bausteine der Materie betrachtet.
	Der Nachweis der Kathodenstrahlteilchen zeigte jedoch, dass Atome eine innere
	Struktur besitzen müssen.
	
	Das Elektron war damit das erste klar identifizierte subatomare Teilchen. Die
	klassische Vorstellung vom Atom als letztem Baustein der Natur war nicht mehr haltbar.
\end{HistoryBox}